\chapter{Introduction}
\label{chap:introduction}

\section{Nightmare Disorder and Imagery Rehearsal Therapy}
\label{sec:intro-irt}
% ~5% prevalence, comorbidity with PTSD/depression
% IRT as gold-standard: structured protocol with rescripting at its core
% Clinician shortage -> scalable AI interventions as potential solution

\section{AI-Assisted Psychotherapy}
\label{sec:intro-ai-therapy}
% LLM-based therapeutic delivery (reDreamAI project as applied context)
% Distinction: general-purpose conversational AI vs. protocol-driven therapeutic agents
% Deployment in clinical settings introduces regulatory constraints (EU AI Act, data sovereignty)

\section{The Evaluation Gap}
\label{sec:intro-eval-gap}
% Standard NLP benchmarks do not assess therapeutic reliability
% Missing: intent stability under stochastic variation, protocol adherence, consistent persona
% Need for domain-specific in silico clinical trials

\section{Research Objectives}
\label{sec:intro-objectives}
% RQ: How consistent is LLM clinical reasoning across stochastic runs in a structured therapeutic protocol?
% Contribution 1: Three-level hierarchical evaluation framework for therapeutic AI stability
% Contribution 2: Quantitative comparison across model classes (sovereign, proprietary, open-weight)
% Scope: isolated rescripting phase (cognitive core of IRT)
